\chapter{总结与展望}

\section{总结}

本文针对移动机器人轨迹跟踪控制中的模型失配问题，提出了一种基于在线稀疏高斯过程回归的模型预测控制框架（GP-MPC），
实现了数据驱动的残差动力学学习与模型预测控制的深度集成。全文的主要工作与贡献总结如下。

\textbf{第二章}建立了差速驱动移动机器人在Frenet坐标系下的运动学模型，推导了其线性化离散状态空间表示，
并介绍了模型预测控制的基本框架，包括滚动优化机制、代价函数设计、约束处理方法以及实时迭代求解策略，
为后续算法设计奠定了理论基础。

\textbf{第三章}提出了面向在线学习的稀疏高斯过程回归算法。以变分推断为理论框架，
利用诱导点处的变分均值与协方差作为模型的充分统计量，将全量数据推断转化为低维变分参数的递推优化问题；
设计了支持增量更新的变分参数递推机制，使模型能够随新观测数据持续修正后验估计；
进一步提出诱导点的增量扩充策略，在不重新训练的前提下提升模型对新状态区域的表达能力。

\textbf{第四章}在上述学习模块的基础上，构建了完整的学习型轨迹跟踪 GP-MPC 框架。
以"名义模型 $+$ GP残差补偿"的形式重构预测模型，推导了随机输入下GP均值与方差的泰勒近似传播公式，
分别提出了零阶与一阶两种GP近似方案；引入方差感知代价函数，使优化器在执行轨迹跟踪任务时
自动朝模型预测不确定性较低的方向求解；并对机会约束框架下的鲁棒约束收紧方法进行了讨论。

\textbf{第五章}在仿真与实物平台上对所提算法进行了系统验证。实验结果表明：
GP-MPC 在轨迹跟踪精度上显著优于纯名义模型 MPC，验证了在线残差学习对模型失配补偿的有效性；
方差感知代价函数的引入进一步提升了跟踪性能与控制稳定性；
一阶 GP 近似相较于零阶近似在预测精度上有所提升，但二者在控制精度上的差异较小，
表明零阶近似在计算效率与控制性能之间具有较好的折中；
实物实验结果与仿真结论基本吻合，验证了算法在真实物理平台上的适用性。

\section{展望}

尽管本文所提 GP-MPC 框架在仿真与实物实验中均取得了良好效果，但仍存在若干值得深入研究的问题，
主要体现在以下三个方面。

\textbf{闭环稳定性的理论分析。}
本文在机会约束一节中对GP-MPC框架的稳定性进行了初步讨论，
但受限于在线学习场景下GP后验的时变性，严格的闭环稳定性证明尚未完成。
后续工作可考虑在有界GP预测误差的假设下，
借助输入到状态稳定性（ISS）或终端约束集方法，
对本文框架的递归可行性与均方稳定性给出更为严格的理论保证。

\textbf{超参数的自适应在线优化。}
本文中GP模型的核函数超参数（长度尺度、信号方差、噪声方差）在离线阶段通过最大似然估计确定，
并在在线运行阶段保持固定。然而，真实系统的动力学特性可能随工况、负载或环境条件的变化而漂移，
固定超参数难以保证模型在全时域内的预测精度。
未来可探索将超参数优化纳入在线学习循环，
结合变分期望最大化（Variational EM）等方法实现超参数的自适应调整，
以提升GP模型对非平稳动力学的适应能力。这在实时系统上往往需要大量算力，因此在算法层面需要有新的突破使得超参数的在线优化
存在可能性。

\textbf{更复杂场景与更高速度下的扩展验证。}
本文的实验主要在中低速、结构化环境下开展。
在更高速度下，车辆动力学的非线性效应更为显著，控制周期内的计算时间约束也更为苛刻，
GP在线更新与MPC求解的实时性面临更大挑战。
此外，在非结构化环境、动态障碍物或多模态工况下，
残差动力学的分布可能发生显著变化，单一GP模型的表达能力可能受限。
后续工作可考虑引入混合专家GP（Mixture of Experts GP）或多模型切换机制，
并在更高速度和更复杂场景下对算法的鲁棒性与实时性进行系统评估。